\documentclass[aspectratio=169,11pt]{beamer}

% Compile with XeLaTeX (Overleaf: Menu > Compiler > XeLaTeX).
% Keep this file beside the assets folder.
\usepackage{fontspec}
\setmainfont{TeX Gyre Heros}
\setsansfont{TeX Gyre Heros}
\usepackage{graphicx}
\usepackage{booktabs,tabularx}
\usepackage{tikz}
\usetikzlibrary{arrows.meta,positioning,shapes.geometric,fit,calc}
\usepackage{hyperref}
\hypersetup{colorlinks=true,urlcolor=travelgreen}

\definecolor{travelgreen}{HTML}{14685E}
\definecolor{travelgold}{HTML}{F4A11A}
\definecolor{softgreen}{HTML}{EAF5EF}
\definecolor{softgold}{HTML}{FFF5DE}
\definecolor{ink}{HTML}{253C3A}
\definecolor{muted}{HTML}{657B78}

\usetheme{Madrid}
\usecolortheme{default}
\setbeamercolor{normal text}{fg=ink,bg=white}
\setbeamercolor{structure}{fg=travelgreen}
\setbeamercolor{frametitle}{fg=travelgreen,bg=white}
\setbeamercolor{title}{fg=travelgreen}
\setbeamercolor{subtitle}{fg=muted}
\setbeamercolor{block title}{fg=white,bg=travelgreen}
\setbeamercolor{block body}{fg=ink,bg=softgreen}
\setbeamercolor{alerted text}{fg=travelgold}
\setbeamertemplate{navigation symbols}{}
\setbeamertemplate{headline}{}
\setbeamertemplate{footline}{
  \leavevmode
  \hbox{
    \begin{beamercolorbox}[wd=.78\paperwidth,ht=2.6ex,dp=1.1ex,leftskip=1em]{author in head/foot}
      \color{muted}\scriptsize Travel Path \textbullet\ Đồ án Trí tuệ nhân tạo
    \end{beamercolorbox}
    \begin{beamercolorbox}[wd=.22\paperwidth,ht=2.6ex,dp=1.1ex,rightskip=1em plus1fil]{date in head/foot}
      \color{travelgreen}\scriptsize \insertframenumber/20
    \end{beamercolorbox}
  }
}
\setbeamertemplate{itemize item}{\color{travelgold}\small$\blacktriangleright$}
\setbeamertemplate{itemize subitem}{\color{travelgreen}\scriptsize$\bullet$}
\setlength{\parskip}{3pt}
\setlength{\leftmargini}{1.2em}
\setbeamersize{text margin left=0.8cm,text margin right=0.8cm}

\tikzset{
  box/.style={draw=travelgreen,rounded corners=5pt,fill=softgreen,align=center,minimum height=0.82cm,text width=2.8cm},
  goldbox/.style={draw=travelgold,rounded corners=5pt,fill=softgold,align=center,minimum height=0.82cm,text width=2.8cm},
  db/.style={draw=travelgreen,ellipse,fill=softgold,align=center,minimum height=0.75cm,text width=2.7cm},
  arr/.style={-{Latex[length=2.5mm]},thick,draw=travelgreen},
  dashedarr/.style={-{Latex[length=2.3mm]},thick,dashed,draw=muted}
}
\newcommand{\asset}[2][]{\includegraphics[#1]{assets/#2}}
\newcommand{\key}[1]{\textcolor{travelgreen}{\textbf{#1}}}

\title[Travel Path]{XÂY DỰNG HỆ THỐNG TRỢ LÝ DU LỊCH\\THÔNG MINH VỚI AI}
\subtitle{Báo cáo thuyết trình đồ án Trí tuệ nhân tạo}
\author{Bùi Thị Quỳnh Như \and Cao Trúc Ny}
\institute{Trường Đại học Kiến trúc Đà Nẵng\\Khoa Công nghệ thông tin -- Lớp 23CT4}
\date{Đà Nẵng -- 2026}

\begin{document}

% 1
\begin{frame}[plain]
  \begin{tikzpicture}[remember picture,overlay]
    \fill[softgreen] (current page.south west) rectangle (current page.north east);
    \fill[travelgreen] ([xshift=9.3cm]current page.south west) rectangle (current page.north east);
    \fill[travelgold] ([xshift=9.15cm]current page.south west) rectangle ([xshift=9.42cm]current page.north east);
  \end{tikzpicture}
  \vspace{0.3cm}
  \begin{columns}[T,onlytextwidth]
    \column{0.62\textwidth}
    {\color{travelgreen}\small\bfseries ĐỒ ÁN TRÍ TUỆ NHÂN TẠO}\\[0.45cm]
    {\color{ink}\LARGE\bfseries XÂY DỰNG HỆ THỐNG\\TRỢ LÝ DU LỊCH\\THÔNG MINH VỚI AI}\\[0.4cm]
    {\large\color{muted}Travel Path -- Lên lịch trình theo nhu cầu,\\sở thích và ngân sách}
    \vfill
    \small
    \textbf{Giảng viên hướng dẫn:} TS. Đỗ Phúc Hào\\
    \textbf{Sinh viên:} Bùi Thị Quỳnh Như -- 2351220202\\
    \hspace{1.66cm}Cao Trúc Ny -- 2251220201
    \column{0.34\textwidth}
    \centering
    \vspace{0.5cm}
    \asset[width=0.67\linewidth]{logo_truong.png}\\[0.5cm]
    {\color{white}\Huge\bfseries TRAVEL\\PATH}\\[0.2cm]
    {\color{white}\large Trợ lý du lịch của bạn}
  \end{columns}
\end{frame}

% 2
\begin{frame}{Nội dung trình bày}
  \framesubtitle{Ba phần chính của đồ án}
  \begin{columns}[T,onlytextwidth]
    \column{.32\textwidth}
      \begin{block}{01 · Tổng quan}
        Bối cảnh, bài toán, mục tiêu và phạm vi.
      \end{block}
    \column{.32\textwidth}
      \begin{block}{02 · Phương pháp}
        Kiến trúc, dữ liệu, AI hội thoại, gợi ý và tính chi phí.
      \end{block}
    \column{.32\textwidth}
      \begin{block}{03 · Kết quả}
        Giao diện, kiểm chứng, giới hạn và hướng phát triển.
      \end{block}
  \end{columns}
  \vspace{0.5cm}
  \begin{center}
    \Large\color{travelgreen}\bfseries
    Từ câu chat tự nhiên đến một lịch trình có thể xem,\\điều chỉnh và lưu theo tài khoản.
  \end{center}
\end{frame}

% 3
\begin{frame}{Bối cảnh và lý do chọn đề tài}
  \begin{columns}[T,onlytextwidth]
    \column{.57\textwidth}
      \begin{itemize}
        \item Lập kế hoạch du lịch đòi hỏi phối hợp điểm đến, thời gian, lưu trú, ăn uống và chi phí.
        \item Thông tin thường nằm ở nhiều nguồn; người dùng phải tự đối chiếu và sắp xếp.
        \item Khi đổi ngân sách, số ngày hoặc sở thích, lịch trình cần được cập nhật đồng bộ.
      \end{itemize}
      \vspace{0.2cm}
      \begin{alertblock}{Lý do chọn đề tài}
        Bài toán kết hợp xử lý ngôn ngữ tự nhiên, hệ gợi ý, dữ liệu địa điểm và tối ưu lịch trình.
      \end{alertblock}
    \column{.39\textwidth}
      \begin{tikzpicture}[node distance=0.32cm]
        \node[box,text width=3.7cm] (a) {Nhiều nguồn thông tin};
        \node[box,below=of a,text width=3.7cm] (b) {Nhiều ràng buộc chuyến đi};
        \node[goldbox,below=of b,text width=3.7cm] (c) {Cần một quy trình lập kế hoạch thống nhất};
        \draw[arr] (a)--(b);
        \draw[arr] (b)--(c);
      \end{tikzpicture}
  \end{columns}
  \vspace{0.25cm}
  {\scriptsize\color{muted}Căn cứ bối cảnh được trích dẫn trong báo cáo: Chiến lược phát triển du lịch Việt Nam đến năm 2030; thống kê du lịch; Luật Du lịch 2017.}
\end{frame}

% 4
\begin{frame}{Bài toán và mục tiêu}
  \begin{block}{Đầu vào}
    \centering Điểm đến $d$ \quad Số ngày $n$ \quad Số người $p$ \quad Ngân sách $b$ \quad Sở thích $S$
  \end{block}
  \begin{center}
    \begin{tikzpicture}[node distance=0.5cm]
      \node[box,text width=2.45cm] (i) {Yêu cầu chuyến đi};
      \node[box,right=of i,text width=2.45cm] (x) {Diễn giải và tra cứu};
      \node[goldbox,right=of x,text width=2.45cm] (o) {Lịch trình + chi phí};
      \draw[arr] (i)--(x);
      \draw[arr] (x)--(o);
    \end{tikzpicture}
  \end{center}
  \begin{itemize}
    \item Tạo lịch trình theo ngày, khung giờ và sở thích.
    \item Ước lượng chi phí từ dữ liệu sẵn có; chỉ rõ khoản thiếu hoặc vượt ngân sách.
    \item Cho phép điều chỉnh, lưu lịch sử riêng cho tài khoản và xuất PDF.
  \end{itemize}
\end{frame}

% 5
\begin{frame}{Phạm vi chức năng}
  \begin{columns}[T,onlytextwidth]
    \column{.49\textwidth}
      \begin{block}{Phía người dùng}
        \begin{itemize}
          \item Đăng ký OTP email, đăng nhập, khôi phục mật khẩu.
          \item Tạo và điều chỉnh plan bằng hội thoại.
          \item Xem chi phí, lưu lịch sử, tải PDF, xóa plan.
        \end{itemize}
      \end{block}
    \column{.49\textwidth}
      \begin{block}{Phía hệ thống}
        \begin{itemize}
          \item Tra cứu dữ liệu địa điểm và giá.
          \item Xếp lịch theo ngày, giờ hoạt động và khoảng cách.
          \item Phân quyền người dùng và quản trị viên.
        \end{itemize}
      \end{block}
  \end{columns}
  \begin{alertblock}{Giới hạn cần nêu}
    Giá và địa chỉ phụ thuộc dữ liệu nguồn; lịch trình là gợi ý, cần người dùng xác nhận trước chuyến đi.
  \end{alertblock}
\end{frame}

% 6
\begin{frame}{Tác nhân và các ca sử dụng}
  \centering
  \begin{tikzpicture}[node distance=0.45cm and 0.8cm]
    \node[box,text width=2.5cm] (user) {Người dùng\\\scriptsize chat · tạo plan · lưu · tải PDF};
    \node[box,right=of user,yshift=0.75cm,text width=2.7cm] (auth) {Tài khoản\\\scriptsize đăng ký · OTP · đăng nhập};
    \node[box,right=of user,yshift=-0.75cm,text width=2.7cm] (plan) {Lịch trình\\\scriptsize tạo · chỉnh sửa · xóa · tải};
    \node[goldbox,right=of plan,yshift=1.5cm,text width=2.5cm] (admin) {Quản trị viên\\\scriptsize xem người dùng};
    \node[db,right=of plan,yshift=-1.2cm,text width=2.7cm] (data) {Dữ liệu địa điểm\\và lịch sử theo tài khoản};
    \draw[arr] (user)--(auth);
    \draw[arr] (user)--(plan);
    \draw[arr] (admin)--(data);
    \draw[arr] (plan)--(data);
  \end{tikzpicture}
  \vspace{0.35cm}
  \begin{center}\small
    Phân tách quyền giúp mỗi tài khoản chỉ xem và thao tác trên lịch sử của chính mình.
  \end{center}
\end{frame}

% 7
\begin{frame}{Kiến trúc tổng thể}
  \centering
  \begin{tikzpicture}[node distance=0.5cm and 0.72cm]
    \node[box,text width=2.5cm] (web) {Frontend\\HTML · CSS · JavaScript};
    \node[box,right=of web,text width=2.7cm] (api) {Backend API\\Python · FastAPI};
    \node[box,above right=0.3cm and 0.75cm of api,text width=2.5cm] (ai) {Ollama\\diễn giải câu chat};
    \node[db,below right=0.3cm and 0.75cm of api,text width=2.8cm] (db) {PostgreSQL / PostGIS\\tài khoản · địa điểm · plan};
    \node[goldbox,right=of db,yshift=1.2cm,text width=2.5cm] (route) {OSRM tùy chọn\\tuyến đường};
    \draw[arr] (web)--node[above,font=\scriptsize]{HTTP / JSON}(api);
    \draw[arr] (api)--(ai);
    \draw[arr] (api)--(db);
    \draw[dashedarr] (api)--(route);
  \end{tikzpicture}
  \vspace{0.3cm}
  \begin{itemize}
    \item Backend điều phối nghiệp vụ, xác thực quyền và truy xuất dữ liệu.
    \item AI hỗ trợ hiểu câu chat; địa điểm và giá phải lấy từ dữ liệu hệ thống.
    \item Khi dịch vụ AI hoặc định tuyến không sẵn sàng, hệ thống có nhánh xử lý dự phòng.
  \end{itemize}
\end{frame}

% 8
\begin{frame}{Công nghệ sử dụng}
  \small
  \begin{tabularx}{\textwidth}{>{\bfseries}p{3.2cm}p{3cm}X}
    \toprule
    Công nghệ & Vai trò & Lý do dùng trong đồ án\\
    \midrule
    HTML, CSS, JavaScript & Giao diện web & Xây dựng chat, trang kế hoạch, tài khoản và lịch sử.\\
    Python, FastAPI & Backend / API & Tổ chức nghiệp vụ và cung cấp các tuyến API.\\
    PostgreSQL, PostGIS & Cơ sở dữ liệu & Lưu tài khoản, địa điểm, vị trí và dữ liệu plan.\\
    Ollama & Mô hình ngôn ngữ cục bộ & Hỗ trợ diễn giải yêu cầu tiếng Việt; có nhánh dự phòng.\\
    Docker Compose & Môi trường chạy & Khởi chạy các dịch vụ phụ thuộc theo cấu hình dự án.\\
    \bottomrule
  \end{tabularx}
  \vspace{0.35cm}
  \begin{alertblock}{Nguyên tắc thiết kế}
    Chia tách giao diện, nghiệp vụ, dữ liệu và AI để dễ thay thế, cấu hình và kiểm tra từng phần.
  \end{alertblock}
\end{frame}

% 9
\begin{frame}{Mô hình dữ liệu khái quát}
  \centering
  \begin{tikzpicture}[node distance=0.7cm and 1.2cm]
    \node[box,text width=2.7cm] (users) {\textbf{users}\\\scriptsize id · email · password\_hash · role};
    \node[box,right=of users,text width=2.7cm] (plans) {\textbf{saved\_travel\_plans}\\\scriptsize id · user\_id · destination · plan};
    \node[goldbox,below=of users,text width=2.7cm] (places) {\textbf{places}\\\scriptsize tên · địa chỉ · loại · giá};
    \node[goldbox,right=of places,text width=2.7cm] (prices) {\textbf{price\_sources}\\\scriptsize giá trị · nguồn · thời điểm};
    \draw[arr] (users)--node[above,font=\scriptsize]{1 : N}(plans);
    \draw[arr] (places)--(prices);
    \draw[dashedarr] (plans)--(places);
  \end{tikzpicture}
  \vspace{0.35cm}
  \begin{itemize}
    \item \texttt{user\_id} liên kết plan với chủ tài khoản.
    \item Giá và vị trí là thuộc tính dữ liệu cần kiểm tra; thiếu giá không được hiểu thành giá bằng 0.
    \item Sơ đồ trình bày các thực thể khái quát, không thay thế ERD đầy đủ của toàn bộ schema.
  \end{itemize}
\end{frame}

% 10
\begin{frame}{Luồng xử lý yêu cầu hội thoại}
  \centering
  \begin{tikzpicture}[node distance=0.36cm]
    \node[box,text width=2.35cm] (a) {Câu chat tiếng Việt};
    \node[box,right=of a,text width=2.35cm] (b) {Trích xuất\\đích · ngày · người · ngân sách};
    \node[box,right=of b,text width=2.35cm] (c) {Kiểm tra đủ\\thông tin};
    \node[goldbox,right=of c,text width=2.35cm] (d) {Tra cứu dữ liệu\\và lập lịch};
    \node[box,below=0.45cm of c,text width=2.65cm] (e) {Hỏi bổ sung nếu còn thiếu};
    \draw[arr] (a)--(b);
    \draw[arr] (b)--(c);
    \draw[arr] (c)--node[above,font=\scriptsize]{đủ}(d);
    \draw[dashedarr] (c)--node[right,font=\scriptsize]{chưa đủ}(e);
    \draw[dashedarr] (e.west) -| ([xshift=-0.1cm]a.south) -- (a.south);
  \end{tikzpicture}
  \vspace{0.5cm}
  \begin{block}{Ví dụ đầu vào}
    “Đà Lạt 3 ngày, 2 người, ngân sách 5 triệu, thích thiên nhiên và cà phê.”
  \end{block}
  {\small Mô hình ngôn ngữ chỉ giúp diễn giải; kết quả địa điểm được đối chiếu với cơ sở dữ liệu.}
\end{frame}

% 11
\begin{frame}{Gợi ý địa điểm và sắp lịch}
  \begin{columns}[T,onlytextwidth]
    \column{.53\textwidth}
      \begin{enumerate}
        \item Lọc theo tỉnh/thành và loại địa điểm.
        \item Ưu tiên nơi phù hợp sở thích.
        \item Kiểm tra địa chỉ, giá và giờ hoạt động có sẵn.
        \item Gom các điểm gần nhau để hạn chế đi lại.
        \item Phân hoạt động vào các khung giờ trong ngày.
      \end{enumerate}
    \column{.43\textwidth}
      \begin{tikzpicture}[node distance=0.35cm]
        \node[box,text width=3.6cm] (pref) {Sở thích + địa bàn};
        \node[box,below=of pref,text width=3.6cm] (filter) {Lọc dữ liệu hợp lệ};
        \node[goldbox,below=of filter,text width=3.6cm] (rank) {Xếp hạng và gom cụm};
        \node[box,below=of rank,text width=3.6cm] (slot) {Gán ngày và khung giờ};
        \draw[arr] (pref)--(filter);
        \draw[arr] (filter)--(rank);
        \draw[arr] (rank)--(slot);
      \end{tikzpicture}
  \end{columns}
  \vspace{0.2cm}
  \begin{alertblock}{Điểm cần lưu ý}
    Đây là bộ lọc và xếp hạng theo tiêu chí nghiệp vụ; không tuyên bố là mô hình học sâu đã huấn luyện.
  \end{alertblock}
\end{frame}

% 12
\begin{frame}{Ước tính ngân sách}
  \begin{columns}[T,onlytextwidth]
    \column{.54\textwidth}
      \begin{block}{Công thức minh họa}
        \[
        C_{\text{chuyến}} =
        C_{\text{lưu trú}}+
        C_{\text{ăn uống}}+
        C_{\text{tham quan}}+
        C_{\text{di chuyển}}
        \]
        \[
        C_{\text{lưu trú}} =
        P_{\text{đêm}}\times(N_{\text{ngày}}-1)
        \]
      \end{block}
      \small Với khoảng giá, hệ thống dùng mức giữa để ước lượng; giá đơn được dùng trực tiếp. Chi phí khách sạn nhân số đêm, không nhân số người theo quy tắc của đồ án.
    \column{.42\textwidth}
      \begin{block}{Ví dụ số học}
        Khách sạn: 80.000--300.000 đ/đêm\\
        Mức giữa: 190.000 đ/đêm\\
        3 ngày = 2 đêm\\[2pt]
        \textbf{Lưu trú dự kiến: 380.000 đ}
      \end{block}
      \scriptsize Ví dụ minh họa phép tính, không phải báo giá thực tế.
  \end{columns}
  \vspace{0.1cm}
  \begin{center}\small
    Nếu thiếu giá, hiển thị “chưa đủ dữ liệu”; nếu vượt ngân sách, nêu rõ phần chênh lệch.
  \end{center}
\end{frame}

% 13
\begin{frame}{Đăng ký, đăng nhập và OTP}
  \centering
  \begin{tikzpicture}[node distance=0.65cm and 0.8cm]
    \node[box,text width=2.4cm] (user) {Người dùng\\email · mật khẩu};
    \node[box,right=of user,text width=2.4cm] (web) {Giao diện tài khoản};
    \node[box,right=of web,text width=2.4cm] (api) {API xác thực};
    \node[goldbox,right=of api,text width=2.4cm] (mail) {SMTP\\gửi mã OTP};
    \draw[arr] (user)--node[above,font=\scriptsize]{đăng ký}(web);
    \draw[arr] (web)--node[above,font=\scriptsize]{yêu cầu OTP}(api);
    \draw[arr] (api)--node[above,font=\scriptsize]{gửi mã}(mail);
    \draw[dashedarr] (mail.south) -- ++(0,-0.55) -| node[pos=0.25,below,font=\scriptsize]{nhập OTP}(web.south);
    \draw[dashedarr] (web.south) -- ++(0,-0.95) -| node[pos=0.25,below,font=\scriptsize]{xác minh và tạo tài khoản}(api.south);
  \end{tikzpicture}
  \vspace{0.55cm}
  \begin{itemize}
    \item Mật khẩu được băm trước khi lưu; không hiển thị lại trong lịch sử hay giao diện.
    \item OTP có thời hạn, giới hạn số lần thử và phụ thuộc cấu hình SMTP hoạt động.
    \item Backend xác minh chủ sở hữu plan trước khi trả dữ liệu.
  \end{itemize}
\end{frame}

% 14
\begin{frame}{Kết quả giao diện: trò chuyện}
  \begin{columns}[T,onlytextwidth]
    \column{.64\textwidth}
      \centering
      \asset[width=\linewidth,height=0.73\textheight,keepaspectratio]{ui_chat.png}
    \column{.34\textwidth}
      \vspace{0.25cm}
      \begin{block}{Vai trò}
        Tiếp nhận yêu cầu tự do bằng tiếng Việt.
      \end{block}
      \begin{itemize}
        \item Hiển thị hội thoại và phản hồi.
        \item Người dùng có thể bổ sung sở thích hoặc yêu cầu chỉnh plan.
        \item Thao tác gửi tin nhắn thuận tiện trên màn hình nhỏ.
      \end{itemize}
  \end{columns}
\end{frame}

% 15
\begin{frame}{Kết quả giao diện: tài khoản}
  \begin{columns}[T,onlytextwidth]
    \column{.58\textwidth}
      \centering
      \asset[width=\linewidth,height=0.76\textheight,keepaspectratio]{ui_login.png}
    \column{.4\textwidth}
      \begin{block}{Luồng tài khoản}
        Đăng nhập bằng email và mật khẩu; đăng ký và khôi phục mật khẩu dùng xác minh OTP.
      \end{block}
      \begin{itemize}
        \item Trang chính yêu cầu người dùng đăng nhập.
        \item Plan và lịch sử được gắn theo tài khoản.
        \item Quản trị viên có quyền xem thông tin đăng ký theo phân quyền.
      \end{itemize}
  \end{columns}
\end{frame}

% 16
\begin{frame}{Kết quả giao diện: lịch trình và ngân sách}
  \begin{columns}[T,onlytextwidth]
    \column{.65\textwidth}
      \centering
      \asset[width=\linewidth,height=0.74\textheight,keepaspectratio]{ui_plan_budget.png}
    \column{.33\textwidth}
      \begin{block}{Thông tin thể hiện}
        Điểm đến, số ngày, số người, sở thích, chi phí và hoạt động theo ngày.
      \end{block}
      \begin{itemize}
        \item Khoản lưu trú.
        \item Tổng chi phí dự kiến.
        \item Bình quân mỗi người.
        \item Phần còn lại hoặc vượt ngân sách.
      \end{itemize}
  \end{columns}
\end{frame}

% 17
\begin{frame}{Kiểm chứng kỹ thuật}
  \begin{columns}[T,onlytextwidth]
    \column{.48\textwidth}
      \begin{block}{Kiểm thử tự động}
        \centering
        {\Huge\color{travelgreen}\bfseries 37/37}\\
        \small bài kiểm thử đạt trong lần chạy ghi nhận.
      \end{block}
      \begin{itemize}
        \item Bao gồm kiểm tra bộ gợi ý và nghiệp vụ lập lịch.
        \item Kiểm tra hồi quy các quy tắc chi phí và lọc địa điểm.
      \end{itemize}
    \column{.48\textwidth}
      \begin{block}{Kiểm tra luồng API}
        \centering
        {\Large\color{travelgreen}\bfseries Đã kiểm tra}\\
        \small tạo plan, phản hồi chat và tải PDF.
      \end{block}
      \begin{itemize}
        \item PDF có ảnh địa điểm nhúng.
        \item Smoke test ghi nhận dữ liệu địa điểm sẵn sàng trong môi trường Docker.
      \end{itemize}
  \end{columns}
  \vspace{0.15cm}
  {\scriptsize\color{muted}Số liệu là kết quả lần chạy kiểm tra trong môi trường phát triển, không phải đánh giá hiệu năng tải hay thử nghiệm người dùng.}
\end{frame}

% 18
\begin{frame}{Minh họa hành trình sử dụng}
  \begin{center}
    \begin{tikzpicture}[node distance=0.4cm]
      \node[box,text width=2.35cm] (s1) {1. Đăng nhập};
      \node[box,right=of s1,text width=2.35cm] (s2) {2. Gửi yêu cầu chuyến đi};
      \node[box,right=of s2,text width=2.35cm] (s3) {3. Xem plan và ngân sách};
      \node[goldbox,right=of s3,text width=2.35cm] (s4) {4. Lưu hoặc tải PDF};
      \draw[arr] (s1)--(s2);
      \draw[arr] (s2)--(s3);
      \draw[arr] (s3)--(s4);
    \end{tikzpicture}
  \end{center}
  \vspace{0.55cm}
  \begin{block}{Ví dụ yêu cầu}
    “Đà Nẵng 3 ngày, 2 người, ngân sách 5 triệu; ưu tiên thiên nhiên và cà phê.”
  \end{block}
  \begin{itemize}
    \item Nếu thiếu thông tin quan trọng, trợ lý hỏi thêm thay vì tự điền tùy ý.
    \item Người dùng có thể yêu cầu thêm sở thích, đổi nơi lưu trú hoặc tạo plan mới.
    \item Kết quả được lưu trong lịch sử của tài khoản đang đăng nhập.
  \end{itemize}
\end{frame}

% 19
\begin{frame}{Hạn chế và hướng phát triển}
  \begin{columns}[T,onlytextwidth]
    \column{.48\textwidth}
      \begin{block}{Hạn chế hiện tại}
        \begin{itemize}
          \item Độ chính xác phụ thuộc dữ liệu địa điểm, giá và giờ hoạt động.
          \item Ước lượng chi phí không thay thế xác nhận từ nhà cung cấp.
          \item Chất lượng gợi ý chưa được đánh giá bằng khảo sát diện rộng.
          \item Triển khai công khai cần bổ sung HTTPS, SMTP và giám sát.
        \end{itemize}
      \end{block}
    \column{.48\textwidth}
      \begin{block}{Hướng phát triển}
        \begin{itemize}
          \item Chuẩn hóa địa chỉ và cập nhật nguồn giá định kỳ.
          \item Tối ưu lựa chọn khi có ngân sách trần.
          \item Đo mức phù hợp của gợi ý với người dùng.
          \item Tăng độ ổn định cho PDF, kiểm thử và triển khai.
        \end{itemize}
      \end{block}
  \end{columns}
\end{frame}

% 20
\begin{frame}{Kết luận}
  \begin{center}
    {\LARGE\color{travelgreen}\bfseries Travel Path kết nối hội thoại, dữ liệu\\và lập lịch trình du lịch.}
  \end{center}
  \vspace{0.3cm}
  \begin{itemize}
    \item Hệ thống tiếp nhận nhu cầu, tìm dữ liệu phù hợp và tạo lịch theo ngày.
    \item Chi phí được ước lượng theo quy tắc rõ ràng; thiếu dữ liệu được nêu riêng.
    \item Tài khoản, lịch sử plan và PDF hoàn thiện luồng sử dụng cơ bản.
  \end{itemize}
  \vfill
  \begin{block}{Tài liệu nền tảng tiêu biểu}
    \scriptsize
    FastAPI: \url{https://fastapi.tiangolo.com/}\\
    PostgreSQL/PostGIS: \url{https://www.postgresql.org/docs/} \quad
    OWASP Password Storage: \url{https://cheatsheetseries.owasp.org/cheatsheets/Password_Storage_Cheat_Sheet.html}\\
    Ollama API: \url{https://github.com/ollama/ollama/blob/main/docs/api.md}
  \end{block}
  \begin{center}\color{travelgold}\Large\bfseries Xin cảm ơn thầy và các bạn đã lắng nghe!\end{center}
\end{frame}

\end{document}

